% Chapter — IMU measurement model (FR-23). Follows the mandatory FR-29
% six-section template: purpose; assumptions and domain of validity;
% derivation; discretization and implementation; validation evidence;
% references. Written ahead of the implementation (Phase 6 chapters land
% before their modules, per the FR-29 accretion rule): this chapter is the
% binding specification for the IMU sensor module. The implementing change
% adds the chapter-manifest entry and replaces the prose validation plan
% with a \testid{} evidence table. The equation labels eq:imu:dtheta,
% eq:imu:gyro, eq:imu:accel, eq:imu:gm, eq:imu:arw, eq:imu:quant,
% eq:imu:oadev, eq:imu:gmadev, eq:imu:presetmap, and eq:imu:recovery are
% intended to be echoed verbatim in the module source (FR-29 traceability).
\chapter{Sensors: Inertial Measurement Unit}
\label{ch:sensors-imu}

\section{Purpose}
\label{sec:imu:purpose}

This chapter specifies the FR-23 IMU sensor model: a triad of rate gyros
and a triad of accelerometers that emit \emph{accumulated increments} ---
an angular increment $\Delta\vv{\theta}$ and a velocity increment
$\Delta\vv{v}$ per sample interval --- corrupted by a deterministic,
seeded error chain comprising turn-on bias, first-order Gauss--Markov
in-run bias, scale factor, small-angle misalignment, angular/velocity
random walk, and quantization. Increments rather than point-sampled rates
are the industry-standard strapdown interface (Titterton and
Weston~\cite{titterton2004}; Savage~\cite{savage1998a}): the sensor
integrates the true angular rate and specific force across the interval,
so intra-interval motion (coning, sculling) is preserved in the signal
handed to the navigation algorithms of Chapters~\ref{ch:gnc-builtin}
and~\ref{ch:ekf} instead of being aliased away. The chapter also derives
the Allan-variance identification machinery --- the overlapping Allan
deviation, the white-noise and bias-instability signatures, and the
recovery procedure --- that Phase~6 exit criterion~1 gates: a
\qty{e4}{\second} static record must return the configured ARW and
bias-instability coefficients to within $\pm 10\,\%$.

\section{Assumptions and domain of validity}
\label{sec:imu:domain}

Assumptions:
\begin{enumerate}
  \item \textbf{Sample schedule.} The IMU emits one increment pair per GNC
    major cycle (decision D-5): the sample interval is exactly the control
    period, $\Delta t_s = T_c$, and sample boundaries coincide with cycle
    boundaries, on which the integrator is already forced to terminate.
    Faster-than-cycle IMU output is deliberately out of scope for v1.
  \item \textbf{Rigid mount at the center of mass.} The sensor triads sit
    at the composite center of mass with a fixed body-frame orientation.
    The size effect (lever-arm acceleration
    $\vv{\omega}\times(\vv{\omega}\times\vv{r})
    + \dot{\vv{\omega}}\times\vv{r}$ at an off-CG station $\vv{r}$) is
    excluded; its magnitude is bounded by
    $\norm{\vv{\omega}}^2 r + \norm{\dot{\vv{\omega}}} r$, e.g.\
    $\le \qty{2e-2}{\metre\per\second\squared}$ for
    $r = \qty{0.5}{\metre}$ at $\norm{\vv{\omega}} =
    \qty{0.2}{\radian\per\second}$, and a configuration placing the IMU
    far off-CG on a tumbling vehicle is outside this model's domain.
  \item \textbf{Static error coefficients.} Scale factor, misalignment,
    and the noise coefficients are constant over a run: no temperature,
    g-sensitivity, or vibration rectification dependence. These belong to
    a higher-fidelity sensor tier, not v1 (PRD Non-Goals discipline).
  \item \textbf{Gauss--Markov in-run bias.} The in-run bias is a
    first-order Gauss--Markov (exponentially correlated) process, held
    constant across each sample interval (zero-order hold). The hold is
    valid for $\Delta t_s \ll \tau_c$; the per-interval error it commits
    relative to the continuous process is $O(\Delta t_s/\tau_c)$ of the
    bias increment and is negligible for the supported regime
    $\Delta t_s \le \qty{1}{\second}$, $\tau_c \ge \qty{10}{\second}$.
  \item \textbf{No rate random walk.} FR-23 names ARW/VRW, bias states,
    scale factor, misalignment, and quantization; a rate-random-walk term
    is deliberately absent. Consequently the model's Allan deviation has
    no rising right-hand branch and no interior minimum; the
    bias-instability identification of
    Section~\ref{sec:imu:allanderivation} therefore evaluates the Allan
    deviation at the analytically known Gauss--Markov peak abscissa
    rather than searching for a global minimum.
\end{enumerate}

Domain of validity: any trajectory the core can propagate, with
$\Delta t_s \ll \tau_c$ (both gyro and accelerometer correlation times)
and quadrature fidelity per the bound of
equation~\eqref{eq:imu:quaderr}. Out-of-domain configurations
(non-positive $\Delta t_s$, $\tau_c \le 0$, negative noise densities or
quanta) are rejected by the FR-15 configuration validator before the run
starts; inside the deterministic loop the model is a total function that
never throws. All quantities are SI (\unit{\radian},
\unit{\metre\per\second}); customary data-sheet units enter only through
the documented conversions after
equation~\eqref{eq:imu:arw}.

\section{Derivation}
\label{sec:imu:derivation}

\subsection{Truth increments}
\label{sec:imu:truth}

Over the sample interval $[t_{k-1}, t_k]$, $t_k - t_{k-1} = \Delta t_s$,
the ideal sensor accumulates
\begin{equation}
  \Delta\vv{\theta}_k = \int_{t_{k-1}}^{t_k} \vv{\omega}^{B}(\tau)\,
    \mathrm{d}\tau,
  \qquad
  \Delta\vv{v}_k = \int_{t_{k-1}}^{t_k} \vv{f}^{B}(\tau)\,\mathrm{d}\tau,
  \label{eq:imu:dtheta}
\end{equation}
where $\vv{\omega}^{B}$ is the true body rate and $\vv{f}^{B}$ is the
true specific force: the total non-gravitational acceleration of the
center of mass, resolved in body axes,
\begin{equation}
  \vv{f}^{B} = \dcm{I}{B}(\quat{q}_{i2b})\,
    \bigl(\vv{a}^{I} - \vv{g}^{I}(\vv{r})\bigr),
  \label{eq:imu:specificforce}
\end{equation}
with $\vv{a}^{I}$ the total inertial acceleration and $\vv{g}^{I}$ the
gravitational acceleration (central body plus configured third bodies,
Chapters~\ref{ch:gravity} and~\ref{ch:thirdbody}). Thrust, aerodynamic,
and solar-radiation-pressure accelerations are non-gravitational and are
sensed; gravitation is not --- an accelerometer in free fall reads zero.

The integrals in equation~\eqref{eq:imu:dtheta} are what distinguishes an
increment interface from a point-sampled rate: a sample
$\vv{\omega}(t_k)\,\Delta t_s$ discards all intra-interval motion,
whereas the integral carries it. Within the simulation, "truth" is the
numerical trajectory of Chapter~\ref{ch:integrators}, whose accepted
steps tile each sample interval exactly (D-5 forces step termination on
cycle boundaries). The model therefore evaluates
equation~\eqref{eq:imu:dtheta} by trapezoidal accumulation over the
accepted steps: for steps $t_{k-1} = \tau_0 < \tau_1 < \dots < \tau_J =
t_k$ with $h_j = \tau_{j+1} - \tau_j$,
\begin{equation}
  \Delta\vv{\theta}_k \approx \sum_{j=0}^{J-1} \frac{h_j}{2}
    \bigl(\vv{\omega}^{B}(\tau_j) + \vv{\omega}^{B}(\tau_{j+1})\bigr),
  \label{eq:imu:quadrature}
\end{equation}
and identically for $\Delta\vv{v}_k$ with $\vv{f}^{B}$. The trapezoidal
rule's local error gives the per-interval quadrature bound
\begin{equation}
  \norm{\vv{e}_{\mathrm{quad}}}
    \le \frac{\Delta t_s\, h_{\max}^2}{12}\,
      \max_{\tau}\norm{\ddot{\vv{\omega}}^{B}(\tau)},
  \label{eq:imu:quaderr}
\end{equation}
e.g.\ $\le \qty{8.3e-9}{\radian}$ per \qty{0.01}{\second} interval at
$h_{\max} = \qty{0.01}{\second}$ and $\norm{\ddot{\vv{\omega}}} \le
\qty{1}{\radian\per\second\cubed}$ --- far below the noise floor of any
configured sensor grade. In the open-loop attitude modes of
Chapter~\ref{ch:vehicle6dof} the rate is constant across each cycle and
equation~\eqref{eq:imu:quadrature} is exact.

\subsection{Error chain}
\label{sec:imu:errorchain}

The measured gyro increment is, per axis triad and in this order,
\begin{equation}
  \boxed{\;
  \Delta\tilde{\vv{\theta}}_k
    = \mathcal{Q}_{q_\theta}\Bigl[
      (\mat{I}_3 + \mat{M}_g)\,\Delta\vv{\theta}_k
      + \bigl(\vv{b}_{g0} + \vv{b}_{g,k}\bigr)\,\Delta t_s
      + \vv{\eta}_{g,k} \Bigr],
  \;}
  \label{eq:imu:gyro}
\end{equation}
and the measured accelerometer increment is
\begin{equation}
  \Delta\tilde{\vv{v}}_k
    = \mathcal{Q}_{q_v}\Bigl[
      (\mat{I}_3 + \mat{M}_a)\,\Delta\vv{v}_k
      + \bigl(\vv{b}_{a0} + \vv{b}_{a,k}\bigr)\,\Delta t_s
      + \vv{\eta}_{a,k} \Bigr],
  \label{eq:imu:accel}
\end{equation}
with $\mathcal{Q}$ the carry-preserving quantizer of
equation~\eqref{eq:imu:quant}. The terms, in the order they apply:

\paragraph{Scale factor and misalignment.} The combined linear distortion
is
\begin{equation}
  \mat{M} = \mat{S} + \mat{\Gamma},
  \qquad
  \mat{S} = \mathrm{diag}(s_x, s_y, s_z),
  \qquad
  \Gamma_{ii} = 0,\; \Gamma_{ij} = \gamma_{ij}\ (i \ne j),
  \label{eq:imu:mis}
\end{equation}
with dimensionless scale-factor errors $s_i$ and six independent
small-angle misalignment entries $\gamma_{ij}$ (\unit{\radian}). The
general off-diagonal form covers both a rigid rotation of an orthogonal
triad ($\gamma_{ij} = -\gamma_{ji}$, three parameters) and
non-orthogonality of the sensing axes (symmetric part); configurations
supply whichever subset they mean and zero the rest. All entries are
fixed configuration constants, not random draws.

\paragraph{Turn-on bias.} One draw per run, per instrument, per axis, at
initialization:
\begin{equation}
  \vv{b}_{g0} \sim \mathcal{N}\bigl(\vv{0},\,
    \sigma_{g0}^2\,\mat{I}_3\bigr),
  \label{eq:imu:turnon}
\end{equation}
constant for the whole run (units \unit{\radian\per\second}; accelerometer
analog $\sigma_{a0}$ in \unit{\metre\per\second\squared}).

\paragraph{In-run bias (first-order Gauss--Markov).} The continuous
process $\dot{\vv{b}} = -\vv{b}/\tau_c + \vv{w}$ with correlation time
$\tau_c$ and stationary standard deviation $\sigma_{\mathrm{GM}}$ is
discretized \emph{exactly} at the sample interval:
\begin{equation}
  \boxed{\;
  \vv{b}_{g,k} = \varphi\, \vv{b}_{g,k-1} + \vv{w}_k,
  \qquad
  \varphi = e^{-\Delta t_s/\tau_{c,g}},
  \qquad
  \vv{w}_k \sim \mathcal{N}\bigl(\vv{0},\,
    \sigma_{g,\mathrm{GM}}^2 (1 - \varphi^2)\, \mat{I}_3\bigr),
  \;}
  \label{eq:imu:gm}
\end{equation}
initialized from the stationary distribution,
$\vv{b}_{g,0} \sim \mathcal{N}(\vv{0},
\sigma_{g,\mathrm{GM}}^2 \mat{I}_3)$ (one draw per run per axis). This
recursion is the exact solution of the Ornstein--Uhlenbeck stochastic
differential equation sampled at $\Delta t_s$ (Maybeck~\cite{maybeck1979},
ch.~4), so the discrete sequence has exactly the stationary variance
$\sigma_{\mathrm{GM}}^2$ and autocorrelation
$\sigma_{\mathrm{GM}}^2 e^{-|m|\Delta t_s/\tau_c}$ at every
$\Delta t_s$ --- no first-order-hold approximation error. The state
$\vv{b}_{g,k}$ is advanced once per sample, \emph{before} the increment
is formed, and is held across the interval; the bias contribution to the
increment is $(\vv{b}_{g0} + \vv{b}_{g,k})\,\Delta t_s$. The
accelerometer bias follows the same recursion with
$(\sigma_{a,\mathrm{GM}}, \tau_{c,a})$.

\paragraph{Angular and velocity random walk.} The white measurement
noise on each increment is
\begin{equation}
  \boxed{\;
  \vv{\eta}_{g,k} \sim \mathcal{N}\bigl(\vv{0},\,
    \sigma_{\Delta\theta}^2\, \mat{I}_3\bigr),
  \qquad
  \sigma_{\Delta\theta} = N_g \sqrt{\Delta t_s},
  \qquad
  \sigma_{\Delta v} = N_a \sqrt{\Delta t_s},
  \;}
  \label{eq:imu:arw}
\end{equation}
where $N_g$ is the angular-random-walk coefficient in
$\unit{\radian}/\sqrt{\unit{\second}}$ (equivalently
$(\unit{\radian\per\second})/\sqrt{\unit{\hertz}}$, the square root of
the one-sided rate-noise power spectral density) and $N_a$ the
velocity-random-walk coefficient in
$(\unit{\metre\per\second})/\sqrt{\unit{\second}}$. The $\sqrt{\Delta
t_s}$ mapping is the defining property of integrated white noise: the
variance of $\int \eta\,\mathrm{d}t$ grows linearly with the integration
time, so the increment noise standard deviation scales with
$\sqrt{\Delta t_s}$ and the reconstructed-rate noise
$\vv{\eta}/\Delta t_s$ scales with $1/\sqrt{\Delta t_s}$. Data-sheet
conversions, exact by definition:
$1\ \mathrm{deg}/\sqrt{\mathrm{h}} = (\pi/180)/60\
\unit{\radian}/\sqrt{\unit{\second}} \approx 2.9089 \times 10^{-4}\
\unit{\radian}/\sqrt{\unit{\second}}$, and
$1\ \mathrm{deg/h} = (\pi/180)/3600\ \unit{\radian\per\second} \approx
4.8481\times 10^{-6}\ \unit{\radian\per\second}$.

\paragraph{Quantization with residual carry.} The output is rounded to
the quantum $q$ ($q_\theta$ in \unit{\radian}, $q_v$ in
\unit{\metre\per\second}) with the sub-quantum residual carried into the
next sample, componentwise:
\begin{equation}
  \boxed{\;
  s_k = u_k + \rho_{k-1},
  \qquad
  y_k = q \left\lfloor \frac{s_k}{q} + \frac{1}{2} \right\rfloor,
  \qquad
  \rho_k = s_k - y_k,
  \qquad
  \rho_0 = 0,
  \;}
  \label{eq:imu:quant}
\end{equation}
where $u_k$ is the bracketed pre-quantization value in
equation~\eqref{eq:imu:gyro} and $y_k$ the emitted increment. The
rounding rule is round-half-up via the floor function --- a single,
branch-free, deterministic IEEE-754 expression; ties ($s_k/q + 1/2$
integral) round toward $+\infty$ by construction. The residual satisfies
$\rho_k \in (-q/2,\, q/2]$, and the carry makes quantization lossless in
accumulation: $\sum_{k \le K} y_k = \sum_{k \le K} u_k - \rho_K$, so the
accumulated output never drifts from the accumulated input by more than
half a quantum --- exactly how a real rate-integrating output register
behaves. Setting $q = 0$ disables the stage ($y_k = u_k$, no carry
state).

\subsection{Allan variance: definitions and model signatures}
\label{sec:imu:allanderivation}

Exit criterion~1 re-identifies the configured noise coefficients from the
model's own output, so the identification chain is derived here in full.
Let $\theta_n = \sum_{k \le n} \Delta\tilde{\theta}_k$ (one axis) be the
accumulated angle --- the "phase" record, $N_\theta$ points at spacing
$\Delta t_s$, with $\theta_0 = 0$ prepended. For cluster time
$\tau = m\,\Delta t_s$ the \emph{overlapping} Allan variance estimator is
(IEEE Std 952~\cite{ieee952}; Riley~\cite{riley2008})
\begin{equation}
  \hat{\sigma}^2(\tau) = \frac{1}{2\tau^2\,(N_\theta - 2m)}
    \sum_{n=0}^{N_\theta - 2m - 1}
    \bigl(\theta_{n+2m} - 2\,\theta_{n+m} + \theta_n\bigr)^2 ,
  \label{eq:imu:oadev}
\end{equation}
the maximally overlapped form, which reuses every stride-$m$ second
difference and has the smallest estimator variance of the standard
family. Its square root $\hat{\sigma}(\tau)$ is the Allan deviation
(ADEV), with the units of the underlying rate
(\unit{\radian\per\second} here).

\paragraph{White rate noise (ARW).} For white rate noise of coefficient
$N$, the Allan variance is $\sigma^2(\tau) = N^2/\tau$
(IEEE Std 952~\cite{ieee952}), a $-1/2$ log--log ADEV slope, so
\begin{equation}
  \sigma(\tau) = \frac{N}{\sqrt{\tau}},
  \qquad\text{and in particular}\qquad
  N = \sigma(\qty{1}{\second})
  \label{eq:imu:whiteslope}
\end{equation}
when $N$ is expressed in $\unit{\radian}/\sqrt{\unit{\second}}$ --- the
standard read-the-curve-at-one-second identification.

\paragraph{Gauss--Markov bias.} The Allan variance of the first-order
Gauss--Markov rate process is derived directly from its autocorrelation
$R(s) = \sigma_{\mathrm{GM}}^2 e^{-|s|/\tau_c}$. Writing $I_1 =
\int_0^\tau b\,\mathrm{d}t$ and $I_2 = \int_\tau^{2\tau}
b\,\mathrm{d}t$, the Allan variance is $\sigma^2(\tau) =
\mathbb{E}[(I_2 - I_1)^2]/(2\tau^2) = (V - C)/\tau^2$ with $V =
\mathrm{Var}(I_1) = \mathrm{Var}(I_2)$ and $C = \mathrm{Cov}(I_1, I_2)$.
Direct integration of the exponential kernel gives
\begin{equation}
  V = 2\sigma_{\mathrm{GM}}^2\bigl[\tau\,\tau_c
      - \tau_c^2\,(1 - e^{-x})\bigr],
  \qquad
  C = \sigma_{\mathrm{GM}}^2\,\tau_c^2\,(1 - e^{-x})^2,
  \qquad
  x = \tau/\tau_c ,
  \label{eq:imu:gmvc}
\end{equation}
and hence, using $2(1-e^{-x}) + (1-e^{-x})^2 = 3 - 4e^{-x} + e^{-2x}$,
\begin{equation}
  \boxed{\;
  \sigma_{\mathrm{GM}}^2(\tau)
    = \frac{2\,\sigma_{\mathrm{GM}}^2\,\tau_c}{\tau}
      \left[ 1 - \frac{\tau_c}{2\tau}
        \bigl( 3 - 4e^{-\tau/\tau_c} + e^{-2\tau/\tau_c} \bigr)
      \right].
  \;}
  \label{eq:imu:gmadev}
\end{equation}
Asymptotically, $\sigma(\tau) \to \sigma_{\mathrm{GM}}
\sqrt{2\tau/(3\tau_c)}$ for $\tau \ll \tau_c$ (slope $+1/2$) and
$\sigma(\tau) \to \sigma_{\mathrm{GM}} \sqrt{2\tau_c/\tau}$ for $\tau \gg
\tau_c$ (slope $-1/2$). Numerical maximization of
equation~\eqref{eq:imu:gmadev} (verified against a direct simulation of
the recursion~\eqref{eq:imu:gm} during drafting) locates the single
interior peak at
\begin{equation}
  \tau^{*} = 1.8926\,\tau_c,
  \qquad
  \sigma_{\mathrm{GM}}(\tau^{*}) = 0.6174\,\sigma_{\mathrm{GM}} .
  \label{eq:imu:gmpeak}
\end{equation}

\paragraph{Bias instability and the preset mapping.} The data-sheet
bias-instability coefficient $B$ (\unit{\radian\per\second}) is defined
through the classical flicker-noise flat-region relation
(IEEE Std 952~\cite{ieee952})
\begin{equation}
  \sigma_{\min} \approx \sqrt{\frac{2\ln 2}{\pi}}\; B = 0.664\,B ,
  \label{eq:imu:bi}
\end{equation}
the value read at the flat bottom of a measured ADEV curve. This model
realizes in-run bias with a Gauss--Markov process, not flicker noise, so
the convention is anchored at the Gauss--Markov peak: a preset specified
by $(B, \tau_c)$ sets the process strength
\begin{equation}
  \boxed{\;
  \sigma_{\mathrm{GM}} = \frac{0.664282}{0.617364}\, B = 1.0760\, B ,
  \;}
  \label{eq:imu:presetmap}
\end{equation}
so that $\max_\tau \sigma_{\mathrm{GM}}(\tau) = 0.664\,B$ exactly and the
conventional read-out returns the configured coefficient. Because the
model contains no rate random walk
(Section~\ref{sec:imu:domain}), the total ADEV keeps falling beyond
$\tau^{*}$ and has no interior minimum; the flat region exists locally at
$\tau^{*}$, where equation~\eqref{eq:imu:gmadev} is stationary.

\paragraph{Quantization.} With the carry of
equation~\eqref{eq:imu:quant}, the accumulated angle equals the true
accumulated angle plus a bounded residual of variance $q_\theta^2/12$
(uniform over a quantum when dithered by the noise terms), which is white
phase noise; its Allan variance is $\sigma^2(\tau) = 3\,(q_\theta^2/12)/
\tau^2$ (IEEE Std 952~\cite{ieee952}), i.e.
\begin{equation}
  \sigma_{Q}(\tau) = \frac{q_\theta}{2\tau},
  \label{eq:imu:quantadev}
\end{equation}
a $-1$ slope that dominates only at the shortest cluster times.

\subsection{Recovery procedure (exit criterion 1)}
\label{sec:imu:recovery}

The identification test propagates a static truth trajectory
($\vv{\omega}^{B} \equiv \vv{0}$, $\vv{f}^{B} \equiv \vv{0}$) for $T =
\qty{e4}{\second}$ at a sample interval $\Delta t_s$ chosen so that
$m\,\Delta t_s = \qty{1}{\second}$ is representable with integer $m$
(e.g.\ $\Delta t_s = \qty{0.01}{\second}$, $m = 100$), with a fixed seed.
Per axis and per instrument:
\begin{enumerate}
  \item Compute $\hat{\sigma}(\tau)$ from
    equation~\eqref{eq:imu:oadev} on an octave grid $\tau = 2^j
    \Delta t_s$ restricted to $\tau \le T/10$, plus the two anchor
    points $\tau = \qty{1}{\second}$ and $\tau = \tau^{*} =
    1.8926\,\tau_c$ (the configured $\tau_c$ is an input to the
    recovery, which re-identifies the noise \emph{magnitudes}).
  \item Recover the random-walk coefficient at the one-second anchor by
    subtracting the two known model contributions in quadrature:
    \begin{equation}
      \hat{N} = \sqrt{ \hat{\sigma}^2(\qty{1}{\second})
        - \sigma_{\mathrm{GM}}^2(\qty{1}{\second})
        - \sigma_{Q}^2(\qty{1}{\second}) } ,
      \label{eq:imu:recovery}
    \end{equation}
    with $\sigma_{\mathrm{GM}}^2(\cdot)$ from
    equation~\eqref{eq:imu:gmadev} at the configured
    $(\sigma_{\mathrm{GM}}, \tau_c)$ and $\sigma_Q^2(\cdot)$ from
    equation~\eqref{eq:imu:quantadev}.
  \item Recover the bias instability at the peak anchor by subtracting
    the recovered white-noise term and inverting
    equation~\eqref{eq:imu:bi}:
    \begin{equation}
      \hat{B} = \frac{1}{0.664282}
        \sqrt{ \hat{\sigma}^2(\tau^{*}) - \hat{N}^2/\tau^{*}
          - \sigma_{Q}^2(\tau^{*}) } .
      \label{eq:imu:recoverybi}
    \end{equation}
  \item Gate $|\hat{N}/N - 1| \le 0.10$ and $|\hat{B}/B - 1| \le 0.10$.
\end{enumerate}
The estimator dispersion supports the $\pm 10\,\%$ gate: at $\tau =
\qty{1}{\second}$ the record holds $\sim T/(2\tau) = 5000$ independent
second differences (relative ADEV scatter well under $2\,\%$), and at
$\tau^{*}$ with the reference preset $\tau_c = \qty{20}{\second}$
($\tau^{*} = \qty{37.9}{\second}$) about $T/(2\tau^{*}) \approx 132$
(scatter $\approx 6\,\%$). Correlation times much beyond
$\qty{50}{\second}$ push $\tau^{*}$ into the poorly estimated tail of a
\qty{e4}{\second} record and are excluded from the gate scenario (not
from the model). The committed test is seeded and deterministic, so the
gate is bit-stable across reruns (D-10).

\section{Discretization and implementation}
\label{sec:imu:implementation}

Binding implementation notes for the Phase~6 module:
\begin{enumerate}
  \item \textbf{Module and traceability.} The model is a core-side
    \texttt{ISensor} (FR-23) sampled once per major cycle. The
    implementing source echoes the labels \texttt{eq:imu:dtheta},
    \texttt{eq:imu:gyro}, \texttt{eq:imu:accel}, \texttt{eq:imu:gm},
    \texttt{eq:imu:arw}, \texttt{eq:imu:quant} verbatim in comments and
    adds this chapter's manifest entry.
  \item \textbf{Evaluation order per sample, normative.} (1) advance the
    gyro and accelerometer Gauss--Markov states by
    equation~\eqref{eq:imu:gm}; (2) form the bracketed values of
    equations~\eqref{eq:imu:gyro} and~\eqref{eq:imu:accel} from the
    accumulated truth increments; (3) apply the quantizer of
    equation~\eqref{eq:imu:quant}. Truth accumulation
    (equation~\eqref{eq:imu:quadrature}) runs inside the propagation
    loop as steps are accepted; the accumulator resets to zero at each
    sample boundary.
  \item \textbf{Seeded stream and draw schedule, normative.} All draws
    come from the named stream \texttt{sensors.imu} (D-9) through the
    core \texttt{NormalSampler}. Initialization consumes, in order:
    gyro turn-on bias $(x,y,z)$, gyro Gauss--Markov initial state
    $(x,y,z)$, accelerometer turn-on bias $(x,y,z)$, accelerometer
    Gauss--Markov initial state $(x,y,z)$. Each sample consumes, in
    order: gyro Gauss--Markov drive $(x,y,z)$, gyro ARW $(x,y,z)$,
    accelerometer Gauss--Markov drive $(x,y,z)$, accelerometer VRW
    $(x,y,z)$. Draws are unconditional --- a zero coefficient multiplies
    a drawn normal by zero rather than skipping the draw --- so the
    stream schedule is independent of which error terms are enabled and
    every configuration of one consumer leaves all other streams
    untouched.
  \item \textbf{Logged channels.} Measured increments land in the
    reserved \texttt{sensors.*} group (FR-16) with unit-suffixed names;
    the true Gauss--Markov bias states are logged alongside as
    truth-side diagnostics so the Chapter~\ref{ch:ekf} consistency
    evaluation can form full-state estimation errors without access to
    core internals.
  \item \textbf{Determinism.} Fixed evaluation order, no allocation in
    the sample path, one carry state per axis per instrument
    (equation~\eqref{eq:imu:quant}) --- the only sensor state besides
    the Gauss--Markov biases; identical seeds reproduce identical
    increment streams bit-for-bit (D-10, exit criterion~1's
    bit-identity clause).
\end{enumerate}

\section{Validation evidence}
\label{sec:imu:validation}

This chapter precedes its module; the validation-evidence table with
machine-checked test identifiers is added by the implementing change.
The planned gates, mirroring Phase~6 exit criterion~1:
\begin{itemize}
  \item \textbf{Golden increments.} Deterministic error chain
    (all noise coefficients zero; fixed $\mat{M}$, biases, quantum)
    against hand-computed references at committed states, exercising
    every term of equations~\eqref{eq:imu:gyro}--\eqref{eq:imu:quant}
    including the quantizer's carry and tie behavior.
  \item \textbf{Gauss--Markov statistics.} Sample variance and lag-one
    autocorrelation of the recursion~\eqref{eq:imu:gm} over a long
    seeded record against $\sigma_{\mathrm{GM}}^2$ and $\varphi$ within
    two-sided 95\,\% bounds.
  \item \textbf{Allan recovery (exit criterion 1).} The
    Section~\ref{sec:imu:recovery} procedure on a \qty{e4}{\second}
    static record recovers $N_g$, $N_a$, $B_g$, $B_a$ within
    $\pm 10\,\%$; the analytic curve
    $\sqrt{N^2/\tau + \sigma_{\mathrm{GM}}^2(\tau) + \sigma_Q^2(\tau)}$
    overlays the estimated ADEV within estimator scatter across the
    octave grid.
  \item \textbf{Coning preservation.} On a pure coning trajectory
    (Chapter~\ref{ch:rigidbody}, closed form), accumulated increments
    match the analytic integral of the body rate to the quadrature
    bound of equation~\eqref{eq:imu:quaderr}, and are measurably
    distinct from point-sampled $\vv{\omega}(t_k)\Delta t_s$.
  \item \textbf{Bit identity.} Two runs of the same seeded scenario
    produce byte-identical sensor channels (exit criterion~1).
\end{itemize}

\section{References}
\label{sec:imu:references}

The increment interface and strapdown error taxonomy follow Titterton and
Weston~\cite{titterton2004} and Savage~\cite{savage1998a}; the Allan
variance definitions, noise-signature table (white, quantization,
bias-instability relation~\eqref{eq:imu:bi}), and test procedure follow
IEEE Std 952~\cite{ieee952} and Riley~\cite{riley2008}, with the
underlying statistic due to Allan~\cite{allan1966}. The exact
Gauss--Markov discretization~\eqref{eq:imu:gm} is standard
(Maybeck~\cite{maybeck1979}); the Gauss--Markov Allan
variance~\eqref{eq:imu:gmadev} is derived above from the process
autocorrelation. Full bibliographic details appear in the bibliography.
